\documentclass[authoryear,11pt]{elsarticle}

\usepackage{hyperref}

\begin{document}

\begin{frontmatter}

\title{Grasp to Identify: Physical-Only Proprioceptive Recognition
of Fifteen Classes with a LEAP Hand}

\author{G.~Sai~Pranav}
\author{Harish J.~Palanthandalam-Madapusi}

\begin{abstract}
Robotic hands commonly use vision or dedicated tactile sensors to recognise objects. Using 15 household objects, this study tests whether a physical LEAP Hand can distinguish among objects based solely on the joint angles achieved during controlled grasps.  Motor-current thresholds stop an entire finger when resistance is detected.  Three sequential grips on the object provide 48 final angles per trial to a 24-unit hidden-layer classifier. The dataset contains five physical trials for each of 15 labels, comprising 14 objects and an empty-hand condition. In three-fold within-dataset validation, the neural network correctly classified 70 of 75 held-out trials. The decision rule accepted 67 trials as known classes, of which 65 were correct. A nearest-trial baseline obtained 69/75 correct in a leave-one-trial-out evaluation. The class set was selected after preliminary analysis, so these results are internal, post-selection estimates. They do not measure generalisation to a new session, object instance, or unseen class.\end{abstract}

\begin{keyword}
proprioception \sep LEAP Hand \sep household object recognition
\sep physical grasping \sep neural network
\end{keyword}

\end{frontmatter}

\noindent\textbf{Category:} Robotics \& Intelligent Machines

\section{Introduction and Objectives}

A robot can learn about an object through contact, even when a camera
cannot see the contact surfaces. Earlier research has used joint
positions and kinesthetic grasp information for object recognition \citep{Sintov2023, Vasquez2017}
but this project asks a narrower, practical question: are the final
angles after a short, repeatable sequence of physical LEAP Hand grips
enough to distinguish a set of household items? The objective is to
build and evaluate a working physical-only recognizer that uses the
hand's existing joint-position readings. A second objective is to
compare a compact neural classifier with a nearest-recorded-trial
method. The study treats an empty-hand grasp as a class. It predicts
a trained label, not the object's dimensions or un-trained object's identity.

\section{Innovation}

The experiment combines three deliberately different grips,
finger-level stopping on motor-current thresholds, and a classifier
that reads only the 48 final joint angles. This removes any need for
vision or added tactile arrays during recognition, unlike most grasp based recognition found in robotic hands. The system also
retains individual physical trials, supports retaking mistaken
recordings, and reports uncertainty rather than forcing every
measurement into a known class, it has a strict confidence rule, and retakes measurements until it is confident in the identified class.
The innovation is an evaluated physical implementation of this specific sensing and decision pipeline.

\section{Methodology}

We used a four-finger, 16-actuator anthropomorphic dexterous LEAP Hand as created by \citet{ShawAgarwalPathak2023}.
We programmed and ran with the same three
configured grips for each item. Objects were placed near a consistent
palm location. During each grip, all joints advanced toward preset
targets. If a moving joint's measured motor current exceeded its
threshold, the entire finger stopped, while other fingers continued.
The controller then recorded all 16 final angles and moved to the
next grip. Five physical trials were collected for each of 14
household items and for the empty-hand condition, yielding 75 real
48-angle vectors. Motor currents controlled stopping but were not
classifier features.

Each angle was normalized using its configured start and target
values. The neural model had 48 inputs, one 24-unit tanh hidden layer,
and 15 softmax outputs. Training used Adam and small Gaussian angle
perturbations, applied only to training folds. Three-fold validation
kept real trials separate from their synthetic copies. The model's
top class, confidence margin, and distance to recorded examples
determined whether a reading was accepted, uncertain, or unknown.
A nearest-trial baseline was evaluated by leaving out one real trial
at a time. Saved trial CSV files and validation reports supported the
analysis.

\section{Results and Conclusions}

The neural model classified 70 of 75 held-out physical trials
correctly before the uncertainty rule, an internal accuracy of
93.3\%. It accepted 67 readings as known; 65 of these were correct,
while eight readings were uncertain. The separate nearest-trial
evaluation obtained 69 correct labels out of 75. Most corresponding joint angles are close,
while differences occur across multiple joints and grips; this
motivates use of the complete angle vector rather than one
hand-chosen threshold. The results
establish feasibility under these recorded conditions, but
independent sessions, new instances, and unseen classes still
require prospective tests.

\begin{thebibliography}{3}

\bibitem[{Shaw et~al.(2023)Shaw, Agarwal, and Pathak}]{ShawAgarwalPathak2023}
Shaw, K., Agarwal, A., Pathak, D., 2023. Leap hand: Low-cost, efficient, and anthropomorphic hand for robot learning.
\url{https://arxiv.org/abs/2309.06440}

\bibitem[V\'{a}squez and Perdereau(2017)]{Vasquez2017}
V\'{a}squez, A. and Perdereau, V. (2017).
Proprioceptive shape signatures for object manipulation and
recognition purposes in a robotic hand.
\emph{Robotics and Autonomous Systems}, 98, 135--146.
\url{https://doi.org/10.1016/j.robot.2017.06.001}.

\bibitem[Sintov and Meir(2023)]{Sintov2023}
Sintov, A. and Meir, I. (2023).
Simple kinesthetic haptics for object recognition.
\emph{The International Journal of Robotics Research}, 42(7), 537--561.
\url{https://doi.org/10.1177/02783649231182486}.

\end{thebibliography}

\end{document}
